\newcommand{\E}{\mathbb E}


\section{Introduction}

As machine learning is being deployed in increasingly consequential domains (including policing \citep{policing}, criminal sentencing \citep{sentencing}, and lending \citep{credit}), the problem of ensuring that learned models are {\em fair\/} has become urgent.

Approaches to fairness in machine learning can coarsely be divided into two kinds: \emph{statistical} and \emph{individual} notions of fairness. Statistical notions typically fix a
small number of protected demographic groups $\mathcal{G}$ (such as racial groups), and then ask for (approximate)
parity of some statistical measure across all of these groups.
%For binary classification, almost all statistical definitions in the literature can be stated as instantiations of the following definition:
%\begin{definition}
%\label{def:statistical}
%Fix a collection of ``protected groups'' $\mathcal{G} = \{\mathcal{G}_1,\ldots,\mathcal{G}_k\}$. Given $k$ distributions $\mathcal{D}_i$ (one for each group) over feature/label pairs $(x,y) \in \mathcal{X}\times \mathcal{Y}$ and a pair of functions on label pairs $s,g:\mathcal{Y} \times \mathcal{Y}\rightarrow \{0,1\}$, a classifier $A$ is $\epsilon$-statistically fair with respect to $(s,g)$ and $\mathcal{G}$ if for each pair of groups $i,j$:
%$$\E_{(x,y) \sim \mathcal{D}_i}[s(A(x),y)|g(A(x),y) = 1] \leq  \E_{(x,y) \sim \mathcal{D}_j}[s(A(x),y)|g(A(x),y) = 1] + \epsilon$$
%\end{definition}
%\sn{Ok the above seems correct, is this from somewhere? I don't think its very intuitive to have it 5 lines in to the paper}
%\mk{I think starting off with this general definition makes more sense if we have results for calibration as well, so then we have
%three instantiations. Otherwise we could just state the defn for FP and say we also consider stat parity.}
One popular statistical measure asks for equality of false positive or negative rates across all groups in $\mathcal{G}$ (this is also sometimes referred to as an {\em equal opportunity\/} constraint \citep{HPS}). Another asks for equality of classification rates (also known as \emph{statistical parity}).
%Examples of statistical measures definable by functions $s,g$ for which parity might be desired include positive and negative classification rates,  false positive and negative rates, and positive (and negative) predictive value, amongst many others --- see the survey \cite{BerkSurvey} for a more extensive overview.
These statistical notions of fairness are the kinds of fairness definitions most common in the literature (see e.g. \cite{kamiran2012data,hajian2013methodology,KMR16,HPS,FSV16,zafar2017fairness,Chou17}).

One main attraction of statistical definitions of fairness is that they can in principle be obtained and checked without making any assumptions about the underlying population, and hence lead to more immediately actionable algorithmic approaches. On the other hand, individual notions of fairness ask for the algorithm to satisfy some guarantee which binds at the individual, rather than group, level. This often has the semantics that ``individuals who are similar'' should be treated ``similarly'' \citep{DworkFairness}, or ``less qualified individuals should not be favored over more qualified individuals'' \citep{JKMR16}. Individual notions of fairness have attractively strong semantics, but their main drawback is that achieving them seemingly requires more assumptions to be made about the setting under consideration.
%Finding accurate classifiers that satisfy statistical notions of fairness can be computationally challenging\footnote{For example, Woodworth et al. \cite{srebro} show that finding a linear threshold classifier that approximately minimizes hinge loss subject to equalizing false positive rates across populations is computationally hard (assuming that refuting a random $k$-XOR formula is hard).}, but checking that they are satisfied by an algorithm (that is, \emph{auditing} the algorithm) is straightforward --- at least when the set of protected groups $\mathcal{G}$ is small and fixed in advance, as it is in the majority of prior work.

The semantics of statistical notions of fairness would be significantly stronger if they were defined over a large number of {\em subgroups\/},
thus permitting a rich middle ground between fairness only for a small number of coarse pre-defined groups, and
the strong assumptions needed for fairness at the individual level.
Consider the kind of {\em fairness gerrymandering\/} that can occur when we only look for unfairness over a small number of pre-defined groups:
\begin{example}
\label{example}
Imagine a setting with two binary features, corresponding to race (say black and white) and gender (say male and female), both of which are distributed independently and uniformly at random in a population. Consider a classifier that labels an example positive if and only if it corresponds to a black man, or a white woman. Then the classifier will appear to be equitable when one considers either protected attribute alone, in the sense that it labels both men and women as positive 50\% of the time, and labels both black and white individuals as positive 50\% of the time. But if one looks at any \emph{conjunction} of the two attributes (such as black women), then it is apparent that the classifier maximally violates the statistical parity fairness constraint. Similarly, if examples have a binary label that is also distributed uniformly at random, and independently from the features, the classifier will satisfy equal opportunity fairness with respect to either protected attribute alone, even though it maximally violates it with respect to conjunctions of two attributes.
\end{example}

We remark that the issue raised by this toy example is not merely hypothetical.
In our experiments in Section \ref{sec:empirical}, we show that similar violations of fairness on
subgroups of the pre-defined groups can result from the application of standard machine learning
methods applied to real datasets.
To avoid such problems, we would like to be able to satisfy a fairness constraint not just for
the small number of protected groups defined by single protected attributes, but for a combinatorially large or even infinite
collection of structured subgroups definable over protected attributes.

%\mk{Maybe this paragraph is the place to give the toy example where there are two binary features and fairness holds wrt each in
%isolation, but there is a discriminated 2CNF. Then we can forward point to the experimental section and say this can happen on real
%data with even simple classifiers like linear thresholds. Also I might introduce the term ``fairness gerrymandering'' since I like it:)
%and think it's ok as long as it's not in the title.}

In this paper, we consider the problem of {\em auditing\/} binary classifiers for equal opportunity and statistical parity,
and the problem of {\em learning\/}  classifiers subject to these constraints,
when the number of protected groups is large.
There are exponentially many ways of carving up a population into subgroups, and we cannot necessarily identify a small number of these \emph{a priori}
as the only ones we need to be concerned about. At the same time, we cannot insist on any notion of statistical fairness for \emph{every} subgroup of the population: for example, any imperfect classifier could be accused of being unfair to the subgroup of individuals defined ex-post as the set of individuals it misclassified. This simply corresponds to ``overfitting'' a fairness constraint. We note that the individual fairness definition of \cite{JKMR16} (when restricted to the binary classification setting) can be viewed as asking for equalized false positive rates across the singleton subgroups, containing just one individual each\footnote{It also asks for equalized false negative rates, and that the false positive rate is smaller than the true positive rate. Here, the randomness in the ``rates'' is taken entirely over the randomness of the classifier.} --- but naturally, in order to achieve this strong definition of fairness, \cite{JKMR16} have to make structural assumptions about the form of the ground truth. It is, however, sensible to ask for fairness for large \emph{structured} subsets of individuals: so long as these subsets have a bounded VC dimension, the \emph{statistical} problem of learning and auditing fair classifiers is easy, so long as the dataset is sufficiently large. This can be viewed as an interpolation between equal opportunity fairness and the individual ``weakly meritocratic'' fairness definition from \cite{JKMR16}, that does not require making any assumptions about the ground truth.
% \sn{I'm not so sure what we mean by the statistical problem is easy - you mean we have uniform convergence and so for any given group we can estimate the expectation?}
Our investigation focuses on the computational challenges, both in theory and in practice.

\subsection{Our Results}

Briefly, our contributions are:
\begin{itemize}
	\item Formalization of the problem of auditing and learning classifiers	
		for fairness with respect to rich classes of subgroups $\mathcal{G}$.
	\item Results proving (under certain assumptions)
		the computational equivalence of auditing
		$\mathcal{G}$ and (weak) agnostic learning of $\mathcal{G}$.
		While these results imply theoretical intractability of auditing for some
		natural classes $\mathcal{G}$, they also suggest that practical machine
		learning heuristics can be applied to the auditing problem.
	\item Provably convergent algorithms for learning classifiers
		that are fair with respect to $\mathcal{G}$, based on a formulation
		as a two-player zero-sum game between a Learner (the primal player) and
		an Auditor (the dual player). We provide two different algorithms, both of which are based on solving for the equilibrium of this game.
		The first provably converges in a polynomial number of steps and is based on simulation of the game dynamics when the Learner
		uses {\em Follow the Perturbed Leader\/} and the Auditor uses best response; the second
		is only guaranteed to converge asymptotically but is computationally simpler, and involves both players using
		{\em Fictitious Play\/}. % \ar{Edited this paragraph}
	\item An implementation and empirical evaluation of the Fictitious Play algorithm demonstrating its
		effectiveness on a real dataset in which subgroup fairness is a concern.
		
\end{itemize}

In more detail,
we start by studying the computational challenge of simply \emph{checking} whether a given classifier satisfies equal opportunity and statistical parity. Doing this in time linear in the number of protected groups is simple: for each protected group, we need only estimate a single expectation. However, when there are many different protected attributes which can be combined to define the protected groups, their number is combinatorially large\footnote{For example, as discussed in a recent Propublica investigation \citep{propublicafacebook}, Facebook policy protects groups against hate speech if the group is definable as a \emph{conjunction} of protected attributes.  Under the Facebook schema, ``race'' and ``gender'' are both protected attributes, and so the Facebook policy protects ``black women'' as a distinct class, separately from black people and women.  When there are $d$ protected attributes, there are $2^d$ protected groups. As a statistical estimation problem, this is not a large obstacle --- we can estimate $2^d$ expectations to error $\epsilon$ so long as our data set has size $O(d/\epsilon^2)$, but there is now a computational problem.}.

We model the problem by specifying a class of functions $\cG$ defined over a set of $d$ protected attributes. $\cG$ defines a set of protected subgroups.
Each function $g \in \cG$ corresponds to the protected
subgroup $\{x : g_i(x) = 1\}$\footnote{For example, in the case of Facebook's policy, the protected attributes include ``race, sex, gender identity, religious affiliation, national origin, ethnicity, sexual orientation and serious disability/disease'' \citep{propublicafacebook}, and $\cG$ represents the class of boolean conjunctions. In other words, a group defined by individuals having any \emph{subset} of values for the protected attributes is protected.}. The first result of this paper is that for both equal opportunity and statistical parity, the computational problem of \emph{checking} whether a classifier or decision-making algorithm $D$ violates statistical fairness with respect to the set of protected groups $\cG$ is equivalent to the problem of \emph{agnostically learning} $\cG$ \citep{agnostic}, in a strong and
distribution-specific sense. This equivalence has two implications:
\begin{enumerate}
\item First, it allows us to import \emph{computational hardness} results from the learning theory literature. Agnostic learning turns out to be computationally hard in the worst case, even for extremely simple classes of functions $\cG$ (like boolean conjunctions and linear threshold functions). As a result, we can conclude that auditing a classifier $D$ for statistical fairness violations with respect to a class $\cG$ is also computationally hard. This means we should not expect to find a polynomial time algorithm that is always guaranteed to solve the auditing problem.

\item However, in practice, various learning heuristics (like boosting, logistic regression, SVMs, backpropagation for neural networks, etc.) are commonly used to learn accurate classifiers which are known to be hard to learn in the worst case. The equivalence we show between agnostic learning and auditing is \emph{distribution specific} --- that is, if on a particular data set, a heuristic learning algorithm can solve the agnostic learning problem (on an appropriately defined subset of the data), it can be used also to solve the auditing problem on the same data set. %We exploit this connection to provide experimental results, and conduct an audit on several real datasets (Communities and Crime, and the NYPD Stop and Frisk Dataset) for various kinds of statistical unfairness. \ar{Put more details here as the experiments are done}
\end{enumerate}
These results appear in Section \ref{sec:auditing}.

Next, we consider the problem of \emph{learning} a classifier that equalizes false positive or negative
rates across all (possibly infinitely many) sub-groups, defined by a class of functions $\cG$. As per the reductions
described above, this problem is computationally hard in the worst case.

However, under the assumption that we have an efficient oracles which solves the \emph{agnostic learning} problem,
we give and analyze algorithms for this problem based on a game-theoretic formulation.
We first prove that the optimal fair classifier can be found as the equilibrium of a two-player, zero-sum game,
in which the (pure) strategy space of the ``Learner'' player corresponds to classifiers in $\cH$,
and the (pure) strategy space of the ``Auditor'' player corresponds to subgroups defined by $\cG$.
The best response problems for the two players correspond to agnostic learning and auditing, respectively. We show that both problems can be solved with a single call to a \emph{cost sensitive classification oracle}, which is equivalent to an agnostic learning oracle.
We then draw on extant theory for learning in games and no-regret algorithms to derive two different algorithms based on simulating game play in this formulation.
In the first, the Learner employs the well-studied {\em Follow the Perturbed Leader (FTPL)\/} algorithm on an appropriate linearization of
its best-response problem, while the Auditor approximately best-responds to the distribution over classifiers of the Learner at each step.
Since FTPL has a no-regret guarantee, we obtain an algorithm that provably converges in a polynomial number of steps.  % \ar{Edited this paragraph}

While it enjoys strong provable guarantees, this first algorithm is randomized (due to the noise added by FTPL), and the best-response step for the
Auditor is polynomial time but computationally expensive. We thus propose a second algorithm that is deterministic, simpler and faster per step, based on both players adopting the Fictitious Play learning dynamic. This algorithm has weaker theoretical guarantees: it has provable convergence only asymptotically, and not in a polynomial number of steps --- but is more practical and
converges rapidly in practice.
% This means that given a heuristic algorithm capable of solving the agnostic learning problem and the auditing problem, both players can implement the fictitious play dynamic, which requires repeatedly best responding to their opponent's empirical history of play. This dynamic is known to converge to Nash equilibrium in zero-sum games --- and hence, in our setting, to the optimal fair classifier.
The derivation of these algorithms (and their guarantees) appear in Section \ref{sec:learning}. % \ar{Edited this paragraph}



Finally, we implement the Fictitious Play algorithm and demonstrate its practicality by efficiently learning classifiers that approximately equalize false positive rates across any group definable by a linear threshold function on 18 protected attributes in the ``Communities and Crime'' dataset. We use simple, fast regression algorithms as heuristics to implement agnostic learning oracles, and (via our reduction from agnostic learning to auditing) auditing oracles. Our results suggest that it is possible in practice to learn fair classifiers with respect to a large class of subgroups that still achieve non-trivial error. Full details are contained in Section \ref{sec:empirical}, and for a substantially more comprehensive empirical investigation of our method we direct the interested reader to \cite{experimental}.


\subsection{Further Related Work}

Independent of our work, \cite{calibration} also consider a related and complementary notion of fairness that they call ``multicalibration''. In settings in which one wishes to train a real-valued predictor, multicalibration can be considered the ``calibration'' analogue for the definitions of subgroup fairness that we give for false positive rates, false negative rates, and classification rates. For a real-valued predictor, calibration informally requires that for every value $v \in [0,1]$ predicted by an algorithm, the fraction of individuals who truly have a positive label in the subset of individuals on which the algorithm predicted $v$ should be approximately equal to $v$. Multicalibration asks for approximate calibration on every set defined implicitly by some circuit in a set $\cG$. \cite{calibration} give an algorithmic result that is analogous to the one we give for learning subgroup fair classifiers: a polynomial time algorithm for learning a multi-calibrated predictor, given an agnostic learning algorithm for $\cG$. In addition to giving a polynomial-time algorithm, we also give a practical variant of our algorithm (which is however only guaranteed to converge in the limit) that we use to conduct empirical experiments on real data.

Thematically, the most closely related piece of prior work is \cite{heuristicauditing}, who also aim to audit classification algorithms for discrimination in subgroups that have not been pre-defined. Our work differs from theirs in a number of important ways. First, we audit the algorithm for common measures of statistical unfairness, whereas \cite{heuristicauditing} design a new measure compatible with their particular algorithmic technique. Second, we give a formal analysis of our algorithm. Finally, we audit with respect to subgroups defined by a class of functions $\cG$, which we can take to have bounded VC dimension, which allows us to give formal out-of-sample guarantees.  \cite{heuristicauditing} attempt to audit with respect to \emph{all possible} sub-groups, which introduces a severe multiple-hypothesis testing problem, and risks overfitting. Most importantly we give actionable algorithms for learning subgroup fair classifiers, whereas \cite{heuristicauditing} restrict attention to auditing.

Technically, the most closely related piece of work (and from which we take inspiration for our algorithm in Section \ref{sec:learning}) is \cite{reductions}, who show that given access to an agnostic learning oracle for a class $\cH$, there is an efficient algorithm to find the lowest-error distribution over classifiers in $\cH$ subject to equalizing false positive rates across polynomially many subgroups. Their algorithm can be viewed as solving the same zero-sum game that we solve, but in which the ``subgroup'' player plays gradient descent over his pure strategies, one for each sub-group. This ceases to be an efficient or practical algorithm when the number of subgroups is large, as is our case. Our main insight is that an agnostic learning oracle is sufficient to have the both players play ``fictitious play'', and that there is a transformation of the best response problem such that an agnostic learning algorithm is enough to efficiently implement follow the perturbed leader. 

There is also other work showing computational hardness for fair learning problems. Most notably, \cite{srebro} show that finding a linear threshold classifier that approximately minimizes hinge loss subject to equalizing false positive rates across populations is computationally hard (assuming that refuting a random $k$-XOR formula is hard). In contrast, we show that even \emph{checking} whether a classifier satisfies a false positive rate constraint on a particular data set is computationally hard (if the number of subgroups on which fairness is desired is too large to enumerate).

%%% Local Variables:
%%% mode: latex
%%% TeX-master: "main.tex"
%%% End:
